\section{TRIT: orientación ternaria, acarreo y representación cuadrática}
\label{trit:capitulo}

La construcción de APP conserva conjuntamente la posición, las evaluaciones
aditiva y multiplicativa y sus descomposiciones nonádicas. Sobre esos datos
actúa ahora la reducción ternaria orientada. Su cometido es doble: distinguir
las tres clases operativas de una evaluación y conservar la coordenada entera
que permite restituirla. Esta distinción interviene después en la selección
de fase del núcleo topológico de fase, denominado TPK por sus iniciales en
inglés, \emph{Topological Phase Kernel}.

El término \emph{trit} designa una unidad ternaria. En la convención balanceada
de HMT sus valores son \(-1,0,+1\). El mismo alfabeto admite funciones
diferentes según el dominio: lectura de una evaluación aritmética, selección
de un modo operativo o índice de una realización cuadrática. Cada función
se define con su dominio y con la información que conserva.

El desarrollo comienza en las evaluaciones enteras de APP. Las realizaciones
cuadráticas aparecen después como representaciones del tipo trítico ya
determinado. La construcción concreta de sus generadores y la extensión de
Euler se desarrollan en el capítulo~\ref{eul:capitulo}, después de introducir
el calendario y la selección de modos del TPK.

\subsection{Reducción ternaria balanceada de las evaluaciones de APP}
\label{trit:reduccion}

En una posición \((i,j)\), con \(i,j\in\{1,\ldots,9\}\), las evaluaciones
enteras son \(S(i,j)=i+j\) y \(P(i,j)=ij\). Cada una conserva su residuo
positivo módulo nueve y su cociente. La reducción ternaria se aplica a la
evaluación correspondiente y produce otro par de coordenadas.

\begin{definicion}[Descomposición ternaria orientada]
\label{trit:def:balanceada}
Sean
\[
 \mathcal T=\{-1,0,+1\},\qquad
 c_3(n)=\left\lfloor\frac{n+1}{3}\right\rfloor,\qquad
 r_3(n)=n-3c_3(n),\quad n\in\mathbb Z.
\]
La aplicación de descomposición y la de restitución son
\[
 \mathcal B_3:\mathbb Z\longrightarrow\mathcal T\times\mathbb Z,
 \qquad \mathcal B_3(n)=(r_3(n),c_3(n)),
\]
\[
 \mathcal L_3:\mathcal T\times\mathbb Z\longrightarrow\mathbb Z,
 \qquad \mathcal L_3(r,c)=r+3c.
\]
El primer componente es el representante ternario balanceado; el segundo
es el cociente de prolongación.
\end{definicion}

\begin{proposicion}[Restitución exacta y unicidad]
\label{trit:prop:restitucion}
Las aplicaciones \(\mathcal B_3\) y \(\mathcal L_3\) son inversas.
En particular, para todo entero \(n\),
\[
 n=r_3(n)+3c_3(n),\qquad r_3(n)\in\mathcal T,
\]
y esta expresión es única.
\end{proposicion}
\begin{proof}
Si \(c=\lfloor(n+1)/3\rfloor\), entonces
\(3c\le n+1<3c+3\). Por tanto \(-1\le n-3c<2\), y el entero
\(r=n-3c\) pertenece a \(\mathcal T\). Si
\(r+3c=s+3d\), con \(r,s\in\mathcal T\), el entero \(r-s\) es múltiplo
de tres y se encuentra entre \(-2\) y \(2\). Así, \(r=s\) y \(c=d\).
Finalmente, para cualquier par \((r,c)\),
\(\lfloor(r+3c+1)/3\rfloor=c\), pues \(0\le r+1\le2\).
Esto demuestra las dos composiciones inversas.
\end{proof}

La sección orientada del cociente \(\mathbb Z/3\mathbb Z\) queda fijada por
\[
 [0]_3\longmapsto0,\qquad [1]_3\longmapsto+1,
 \qquad [2]_3\longmapsto-1.
\]
La sustitución de \(2\) por \(-1\) va acompañada del cociente necesario:
\(2=-1+3\). La evaluación entera permanece recuperable en el par.

\begin{tablacompleta}
\caption{Descomposición balanceada de las nueve posiciones de una recta APP.}
\label{trit:tabla:nueve}
\begin{tabular}{c*{9}{r}}
\toprule
\(n\)&1&2&3&4&5&6&7&8&9\\
\midrule
\(r_3(n)\)&1&$-1$&0&1&$-1$&0&1&$-1$&0\\
\(c_3(n)\)&0&1&1&1&2&2&2&3&3\\
\bottomrule
\end{tabular}
\notatabla{La coincidencia de representantes ternarios conserva tres
posiciones distintas en cada clase; sus cocientes restituyen los enteros.}
\end{tablacompleta}

Por ejemplo, en \((i,j)=(5,5)\), la suma vale \(10\) y el producto vale
\(25\). Las dos lecturas ternarias dan \(+1\), mientras que sus
descomposiciones completas son
\[
 \mathcal B_3(10)=(1,3),\qquad
 \mathcal B_3(25)=(1,8).
\]
La firma común identifica un tipo residual; la hoja y el cociente distinguen
las dos evaluaciones. El estado de APP conserva además la posición que las
produjo.

\subsubsection{Compatibilidad entre los módulos nueve y tres}

Escribamos la descomposición nonádica de una evaluación como
\[
 n=R_9(n)+9q_9(n),\qquad R_9(n)\in\{1,\ldots,9\}.
\]
Al reducir su representante positivo ternariamente se obtiene
\[
 R_9(n)=r_3(R_9(n))+3c_3(R_9(n)).
\]
Sustituyendo,
\begin{equation}
 r_3(n)=r_3(R_9(n)),\qquad
 c_3(n)=c_3(R_9(n))+3q_9(n).
\label{trit:eq:compatibilidad39}
\end{equation}
La unicidad de la descomposición balanceada demuestra ambas identidades.
La reducción visible módulo tres puede efectuarse sobre el residuo
nonádico, mientras la reconstrucción completa transporta también el
cociente anterior.

En el producto \(5\cdot5=25\), la descomposición nonádica es
\(25=7+9\cdot2\). Como \(7=1+3\cdot2\), la fórmula anterior da
\(25=1+3(2+3\cdot2)=1+3\cdot8\). Los dos órdenes de lectura restituyen
el mismo entero porque los cocientes se componen de manera explícita.

\subsection{Aritmética de los estados ternarios con acarreo}
\label{trit:aritmetica}

La conservación de los pares permite trasladar las operaciones de APP a
las coordenadas balanceadas. El traslado determina el acarreo aditivo y
los términos cruzados del producto.

\begin{definicion}[Acarreo aditivo balanceado]
Para \(r,s\in\mathcal T\), sean
\[
 r\oplus s=r_3(r+s),\qquad
 \chi(r,s)=\frac{r+s-r_3(r+s)}3.
\]
El entero \(\chi(r,s)\) es el acarreo producido al normalizar la suma de
los dos representantes.
\end{definicion}

\begin{tablacompleta}
\caption{Acarreo aditivo de la sección ternaria balanceada.}
\begin{tabular}{c|rrr}
\toprule
\(\chi(r,s)\)&\(s=-1\)&\(s=0\)&\(s=+1\)\\
\midrule
\(r=-1\)&$-1$&0&0\\
\(r=0\)&0&0&0\\
\(r=+1\)&0&0&1\\
\bottomrule
\end{tabular}
\notatabla{Los dos extremos codifican
\(-1-1=1-3\) y \(1+1=-1+3\).}
\end{tablacompleta}

\begin{proposicion}[Operaciones recuperables]
\label{trit:prop:operaciones}
Si \(n=r+3c\) y \(m=s+3d\), con \(r,s\in\mathcal T\), entonces
\begin{align}
 (r,c)\boxplus(s,d)
 &=\bigl(r\oplus s,c+d+\chi(r,s)\bigr),
 \label{trit:eq:suma}\\
 (r,c)\boxtimes(s,d)
 &=\bigl(rs,rd+sc+3cd\bigr)
 \label{trit:eq:producto}
\end{align}
son, respectivamente, las descomposiciones de \(n+m\) y \(nm\).
\end{proposicion}
\begin{proof}
La suma se descompone como
\[
 n+m=r+s+3(c+d)
     =r_3(r+s)+3\bigl(c+d+\chi(r,s)\bigr).
\]
En el producto,
\[
 nm=(r+3c)(s+3d)=rs+3(rd+sc+3cd).
\]
El entero \(rs\) pertenece a \(\mathcal T\). Las dos expresiones están
normalizadas; la unicidad de la proposición~\ref{trit:prop:restitucion}
identifica sus componentes.
\end{proof}

Estas leyes distinguen la acumulación y la multiplicación en el registro.
La suma agrega los cocientes y añade el acarreo residual. El producto
combina también cada representante con el cociente de la otra evaluación.
Ese contenido aritmético se mantiene cuando las operaciones se incorporan
a una transición del TPK.

\begin{ejemplo}[Suma y producto de dos evaluaciones]
De \(2=-1+3\) y \(4=1+3\) resulta
\[
 (-1,1)\boxplus(1,1)=(0,2),\qquad
 (-1,1)\boxtimes(1,1)=(-1,3).
\]
La restitución produce \(6\) y \(8\), respectivamente. En la segunda
operación, el cociente es \((-1)\cdot1+1\cdot1+3\cdot1\cdot1=3\).
\end{ejemplo}

\begin{proposicion}[Identidad de cociclo y asociatividad]
\label{trit:prop:cociclo}
Para \(r,s,t\in\mathcal T\),
\begin{equation}
 \chi(r,s)+\chi(r\oplus s,t)
 =\chi(s,t)+\chi(r,s\oplus t).
\label{trit:eq:cociclo}
\end{equation}
Las operaciones \(\boxplus\) y \(\boxtimes\), junto con sus elementos
neutros e inversos aditivos transportados por \(\mathcal B_3\), hacen
de \(\mathcal T\times\mathbb Z\) un anillo isomorfo a \(\mathbb Z\).
\end{proposicion}
\begin{proof}
Al normalizar \((r+s)+t\), el cociente total es el primer miembro de
\eqref{trit:eq:cociclo}; al normalizar \(r+(s+t)\), es el segundo.
Los representantes finales coinciden porque representan la misma clase
módulo tres. La unicidad de la descomposición completa obliga a que también
coincidan los cocientes. Por la proposición~\ref{trit:prop:operaciones},
\(\mathcal L_3\) es una biyección que conserva suma y producto. Transporta
por ello la asociatividad, distributividad y conmutatividad de las
operaciones enteras. Sus neutros son \((0,0)\) y \((1,0)\).
\end{proof}

La prueba de cociclo expresa una compatibilidad concreta del registro:
la normalización puede efectuarse en etapas sin cambiar la evaluación
final. Si interesa la secuencia de operaciones, el registro conserva
además el orden de las normalizaciones efectuadas.

\subsection{Refinamiento ternario y recuperación a profundidad arbitraria}
\label{trit:refinamiento}

La descomposición se itera sobre el cociente. Dado \(n_0\in\mathbb Z\),
defínanse recursivamente
\[
 r_k=r_3(n_k),\qquad n_{k+1}=c_3(n_k),\qquad k\ge0.
\]
Cada transición satisface \(n_k=r_k+3n_{k+1}\). La composición de estas
igualdades produce una restitución a cualquier profundidad prefijada.

\begin{proposicion}[Restitución de una truncación]
\label{trit:prop:profundidad}
Para todo \(N\ge1\),
\begin{equation}
 n_0=\sum_{k=0}^{N-1}3^kr_k+3^Nn_N.
\label{trit:eq:restitucionN}
\end{equation}
La sucesión \((r_0,\ldots,r_{N-1})\), junto con el cociente terminal
\(n_N\), determina exactamente \(n_0\).
\end{proposicion}
\begin{proof}
Para \(N=1\), la identidad es la descomposición inicial. Si vale para
\(N\), se sustituye \(n_N=r_N+3n_{N+1}\) en el último término y se
obtiene la fórmula para \(N+1\). La restitución reconstruye también los
cocientes intermedios, comenzando en \(n_N\) y aplicando
\(n_k=r_k+3n_{k+1}\) en orden inverso.
\end{proof}

\begin{ejemplo}[Un acarreo que atraviesa tres niveles]
La evaluación \(26\) da sucesivamente
\[
 26=-1+3\cdot9,\qquad 9=0+3\cdot3,
 \qquad 3=0+3\cdot1,\qquad1=1+3\cdot0.
\]
Sus representantes, ordenados de menor a mayor potencia, son
\((-1,0,0,1)\). La restitución es \(26=-1+27\). El símbolo cero de
los niveles intermedios conserva una posición y un cociente determinado.
\end{ejemplo}

\begin{proposicion}[Terminación para un entero fijo]
La iteración del cociente balanceado alcanza cero para todo entero
inicial. Una vez suprimidos los ceros finales, la expansión ternaria
balanceada finita de un entero es única.
\end{proposicion}
\begin{proof}
Para \(|n|\ge2\),
\( |c_3(n)|\le(|n|+1)/3<|n|\).
Para \(n\in\{-1,0,1\}\), el cociente es cero. El valor absoluto
desciende estrictamente mientras es al menos dos, de modo que la
iteración termina. Si dos expansiones representan el mismo entero,
su reducción módulo tres fija el mismo primer representante; al restarlo
y dividir por tres se repite el argumento para todas las posiciones.
\end{proof}

Este resultado se refiere a la aritmética de un entero fijo. En las
prolongaciones del TPK aumenta la longitud del registro y pueden generarse
nuevas evaluaciones. Allí la profundidad corresponde a una sucesión
compatible de estados, cuya construcción requiere las reglas de
supervivencia. La fórmula~\eqref{trit:eq:restitucionN} aporta el mecanismo
local de restitución que esas prolongaciones deben conservar.

\subsection{Dos coordenadas ternarias de una clase nonádica}
\label{trit:nonadico}

La reducción de \(\mathbb Z/9\mathbb Z\) módulo tres tiene tres elementos
en cada fibra. La coordenada adicional puede escribirse como un cociente
ternario residual.

\begin{proposicion}[Coordenadas nonádicas con suma acarreada]
\label{trit:prop:beta}
La aplicación
\[
 \beta:\mathcal T\times\mathbb F_3\longrightarrow\mathbb Z/9\mathbb Z,
 \qquad \beta(r,\bar c)=[r+3c]_9
\]
es una biyección de conjuntos. La suma y el producto nonádicos toman la
forma
\begin{align}
 \beta(r,\bar c)+\beta(s,\bar d)
 &=\beta\bigl(r\oplus s,\overline{c+d+\chi(r,s)}\bigr),
 \label{trit:eq:beta-suma}\\
 \beta(r,\bar c)\beta(s,\bar d)
 &=\beta\bigl(rs,\overline{rd+sc}\bigr).
 \label{trit:eq:beta-producto}
\end{align}
\end{proposicion}
\begin{proof}
Cambiar \(c\) por \(c+3k\) modifica \(r+3c\) en \(9k\), por lo que
la aplicación está bien definida. Si dos imágenes coinciden, la reducción
módulo tres fija el mismo \(r\). La diferencia restante, dividida por
tres, fija la misma clase \(\bar c\). La inyectividad entre dos conjuntos
de nueve elementos prueba la biyección. Las operaciones proceden de
\eqref{trit:eq:suma} y \eqref{trit:eq:producto}; el término \(3cd\) del
cociente multiplicativo se anula al reducirlo módulo tres.
\end{proof}

El acarreo establece la ley de composición entre los dos niveles. Por
ejemplo,
\[
 \beta(1,0)+\beta(1,0)=\beta(-1,1).
\]
La segunda coordenada registra el paso que realiza la suma entera antes
de su lectura residual. En particular, \(\mathbb Z/9\mathbb Z\) tiene
elementos de orden nueve, mientras que el grupo aditivo de
\(\mathbb F_3^2\) tiene exponente tres. La biyección anterior conserva
la suma mediante la ley acarreada explícita.

Las tres posiciones \(3,6,9\) de APP se escriben
\[
 [3]_9=\beta(0,1),\qquad
 [6]_9=\beta(0,2),\qquad
 [9]_9=\beta(0,0).
\]
Sobre el ideal \(3\mathbb Z/9\mathbb Z\), la extracción de la segunda
coordenada es
\[
 [3c]_9\longmapsto[c]_3.
\]
Ésta es una biyección aditiva con \(\mathbb F_3\). Se distingue de la
reducción directa módulo tres, que asigna cero a los tres elementos.
Además, el producto de dos elementos de ese ideal es cero módulo nueve.
Su estructura multiplicativa conserva así una función diferente de la
multiplicación del cuerpo \(\mathbb F_3\).

\subsubsection{La fibra de seis posiciones ternarias}

Una palabra orientada de longitud seis pertenece a \(\mathcal T^6\),
con \(3^6=729\) posibilidades. La elección de representantes induce una
biyección con \(\mathbb F_3^6\). Al agrupar posiciones de dos en dos y
aplicar \(\beta\), se obtiene también una biyección conjuntista con
\((\mathbb Z/9\mathbb Z)^3\). La ley transportada utiliza el acarreo
de cada par.

El censo de palabras efectivamente emitidas por el TPK se realiza después
de construir su emisor y su cronología. La fibra ambiente de \(729\)
elementos fija capacidad coordenada; el conjunto de emisiones y su imagen
ternaria se determinan por las reglas operativas. Esta separación permite
conservar, en el censo ulterior, tanto las multiplicidades como los
registros distintos que comparten una palabra visible.

\subsection{Involución orientada y memoria de transporte}
\label{trit:involucion}

\begin{proposicion}[Compatibilidad con la inversión de signo]
Para todo \(n\in\mathbb Z\),
\[
 \mathcal B_3(-n)=(-r_3(n),-c_3(n)).
\]
\end{proposicion}
\begin{proof}
De \(n=r+3c\) se deduce \(-n=(-r)+3(-c)\). Como \(-r\in\mathcal T\),
la unicidad de la descomposición demuestra la afirmación.
\end{proof}

Para una sucesión de firmas de longitud \(m\), se define
\[
 \mathcal C_{\rm rev}(\tau_1,\ldots,\tau_m)
 =(-\tau_m,\ldots,-\tau_1).
\]
La segunda aplicación restituye el orden y los signos iniciales, de modo
que \(\mathcal C_{\rm rev}^{2}=\operatorname{id}\). En una secuencia
de pares, la misma regla transporta simultáneamente
\((r_k,c_k)\mapsto(-r_{m+1-k},-c_{m+1-k})\). La longitud y cada
evaluación restituible quedan determinadas.

La memoria de división y la memoria de transporte cumplen cometidos
distintos. El par \((r_3(n),c_3(n))\) reconstruye una evaluación entera.
La reconstrucción de una trayectoria necesita además posición inicial,
orientaciones, secuencia de modos y pasos de frontera. Dos recorridos
pueden alcanzar la misma evaluación y conservar registros de transporte
diferentes. El estado enriquecido integra ambos órdenes de información.

En la interpretación temporal adoptada por HMT, \(+1\) corresponde al
retorno con memoria, \(0\) al umbral presente y \(-1\) a la apertura de
prolongaciones conjugadas. Esta asignación fija un lector de régimen.
La evolución efectiva requiere la transición del TPK sobre los registros
completos; sus iteraciones determinan qué fase retorna y qué memoria se
incrementa.

\subsection{Álgebras cuadráticas de los tres regímenes}
\label{trit:cuadraticas}

El índice trítico posee una realización algebraica mediante una relación
cuadrática. Su formulación integral precede a la elección del cuerpo de
escalares y permite separar la información discreta de las realizaciones
analíticas posteriores.

\begin{definicion}[Álgebra cuadrática de tipo trítico]
Para \(\tau\in\mathcal T\), se define
\[
 \mathcal Q_{\tau,\mathbb Z}=\mathbb Z[u]/(u^2+\tau).
\]
Para un anillo conmutativo \(A\), su realización por cambio de escalares es
\[
 \mathcal Q_{\tau,A}=A[u]/(u^2+\tau).
\]
La clase de \(u\), escrita \(J_\tau\), satisface
\(J_\tau^2=-\tau I\).
\end{definicion}

La división por el polinomio mónico \(u^2+\tau\) proporciona una expresión
única \(aI+bJ_\tau\). El producto se calcula reduciendo únicamente el
término cuadrático:
\begin{equation}
 (aI+bJ_\tau)(cI+dJ_\tau)
 =(ac-\tau bd)I+(ad+bc)J_\tau.
\label{trit:eq:producto-cuadratico}
\end{equation}
Esta multiplicación representa el régimen cuadrático. La aritmética de
los pares balanceados conserva, por su parte, los enteros y sus cocientes
mediante \eqref{trit:eq:suma} y \eqref{trit:eq:producto}.

\subsubsection{Realizaciones de nueve elementos}

Sobre \(\mathbb F_3\), las tres álgebras tienen nueve elementos y admiten
la clasificación siguiente:
\[
 \mathcal Q_{+1,\mathbb F_3}\cong\mathbb F_9,\qquad
 \mathcal Q_{0,\mathbb F_3}=\mathbb F_3[\varepsilon]/(\varepsilon^2),
 \qquad
 \mathcal Q_{-1,\mathbb F_3}\cong\mathbb F_3\times\mathbb F_3.
\]

\begin{proposicion}[Distinción de los tres tipos finitos]
\label{trit:prop:tres-tipos}
La primera álgebra es un cuerpo; la segunda contiene un nilpotente no
nulo; la tercera tiene dos idempotentes ortogonales no triviales cuya
suma es la unidad. Las tres son mutuamente no isomorfas.
\end{proposicion}
\begin{proof}
El polinomio \(u^2+1\) toma los valores \(1,2,2\) en \(0,1,2\);
por tanto es irreducible sobre \(\mathbb F_3\). Su cociente es un
cuerpo de grado dos. En el segundo caso, la clase \(\varepsilon\) es
no nula y satisface \(\varepsilon^2=0\). En el tercero,
\(u^2-1=(u-1)(u+1)\), con factores coprimos. La evaluación
\(a+bu\mapsto(a+b,a-b)\) es un isomorfismo a
\(\mathbb F_3\times\mathbb F_3\). Los elementos
\(2(1+u)\) y \(2(1-u)\) son sus dos idempotentes coordenados.
Estas propiedades algebraicas distinguen los tres tipos.
\end{proof}

En el desarrollo de incidencia excepcional, un operador ternario
concretamente construido que satisfaga \(A_W^2=-I\) induce la acción de
\(\mathbb F_9\) sobre su espacio. Si ese espacio es \(\mathbb F_3^6\),
la dimensión sobre \(\mathbb F_9\) es tres: la acción se define por
\((a+b\iota)v=av+bA_Wv\). El operador fija la multiplicación por la
unidad cuadrática. La construcción de Paley correspondiente aporta ese
operador en su lugar causal; la presente clasificación establece el
contenido de la relación que deberá satisfacer.

\subsubsection{Conjugación, norma algebraica e inversa}

En la realización real, la conjugación y la norma algebraica son
\[
 \overline{aI+bJ_\tau}=aI-bJ_\tau,\qquad
 N_\tau(aI+bJ_\tau)=a^2+\tau b^2.
\]

\begin{proposicion}[Norma multiplicativa y restitución inversa]
\label{trit:prop:norma}
Para \(z,w\in\mathcal Q_{\tau,\mathbb R}\),
\[
 z\bar z=N_\tau(z)I,
 \qquad N_\tau(zw)=N_\tau(z)N_\tau(w).
\]
Un elemento \(z\) es invertible exactamente cuando \(N_\tau(z)\ne0\);
en ese caso \(z^{-1}=\bar z/N_\tau(z)\).
\end{proposicion}
\begin{proof}
El producto \eqref{trit:eq:producto-cuadratico} da
\((aI+bJ_\tau)(aI-bJ_\tau)=(a^2+\tau b^2)I\).
La conjugación es multiplicativa y el álgebra es conmutativa, por lo que
\((zw)\overline{zw}=(z\bar z)(w\bar w)\). Si la norma es no nula,
la fórmula exhibe una inversa. Si \(zw=I\), la multiplicatividad da
\(N_\tau(z)N_\tau(w)=1\), que obliga a \(N_\tau(z)\ne0\).
\end{proof}

La representación matricial
\[
 aI+bJ_\tau\longmapsto
 \begin{pmatrix}a&-\tau b\\b&a\end{pmatrix}
\]
es fiel y tiene determinante \(a^2+\tau b^2\). Para \(\tau=+1\), la
forma es positiva definida. Para \(\tau=0\), depende únicamente de
\(a\), mientras \(J_0\) conserva un componente lineal no nulo. Para
\(\tau=-1\), la forma tiene firma \((1,1)\); las direcciones
\(I\pm J_{-1}\) son divisores de cero.

En particular, la palabra \emph{norma} designa aquí una forma cuadrática
multiplicativa. Las normas positivas de espacios de Hilbert, utilizadas
en otras realizaciones del estado, conservan su propia definición.

\subsubsection{Composición de razones entre coordenadas}

Una dirección con componente escalar no nulo se representa por
\(I+uJ_\tau\). La multiplicación produce
\[
 (I+uJ_\tau)(I+vJ_\tau)
 =(1-\tau uv)I+(u+v)J_\tau.
\]
Cuando \(1-\tau uv\ne0\), la razón entre componentes del producto es
\begin{equation}
 u\star_\tau v=\frac{u+v}{1-\tau uv}.
\label{trit:eq:razones}
\end{equation}
El factor escalar \(1-\tau uv\) forma parte del producto completo.
Si se conserva únicamente la razón, ese factor debe registrarse cuando
intervenga en una reconstrucción posterior. Si el denominador se anula,
el producto se examina en sus dos componentes: puede representar una
dirección de componente escalar nulo o el elemento cero.

La clasificación cuadrática prepara así una lectura precisa de los tres
regímenes. La selección de un régimen conserva su índice, la aritmética
conserva sus acarreos y el transporte conserva la secuencia orientada.
El TPK compone esos datos mediante operadores explícitos. Sobre esa
composición se construirán los generadores concretos y la conservación
de formas que articula la extensión de Euler.
